\chapter[Sermon 45. St. John Devours the Book Received at the Angel's Hand, and Prophesies Again to the Gentiles, Nations, and Kings.]{St. John Devours the Book Received at the Angel's Hand, and Prophesies Again to the Gentiles, Nations, and Kings.}
\label{sermon:045}
\markright{Sermon 45. St. John Devours the Book Received at the Angel's ...}

And the voice which I heard from Heaven spoke to me again, and said: Go and take the little book, which is open in the hand of the Angel, who stands upon the sea, and upon the earth. And I went to the Angel and said to him, Give me the little book. And he said to me, Take it, and eat it up, and it shall make your belly bitter, but it shall be in your mouth as sweet as honey. And I took the little book out of the hand of the angel, and did eat it up, and it was in my mouth as sweet as honey. And as soon as I had eaten it, my belly was bitter. And he said to me: You must prophesy again concerning many nations, tongues, and peoples, and to many kings.

\index{Antichrist}\index{John, Saint}\index{Mahomet}This is the third comfort, which is contained in this tenth chapter. For under the person of St. John, it is shown here that the apostolic and evangelical doctrine must be restored in the last times before the judgment against Antichrist and Mahomet. And he might have briefly said, “The apostolic doctrine, as it was preached by John, shall flourish again,” but he preferred to express the same by a beautiful vision, at the last adding a plain and brief exposition of the vision. Which is, “You must preach again, \&c.”

\index{Primasius}\index{Aquinas, Thomas}And all those who expound these things agreeably, first concerning the person of John, who, under Emperor Nero, returned from exile to Asia after about five years, and again preached the gospel. For he lived until the third or fourth year of the reign of Emperor Trajan. Secondly, concerning all preachers before the final judgment, endowed with the spirit and doctrine of Saint John, and constantly professing Christ against Antichrist. Primasius, expounding this passage, says that the certain meaning is directed to Saint John, who must yet be released from exile, not only to bring this revelation to the knowledge of Christ’s church, but also to preach the gospel more deeply to people and nations, to tongues and many kings. Nevertheless, no one doubts that this voice also agrees to the whole Church, which never ought to cease from preaching, and so forth. Thus he says. The ordinary gloss expounds these words: although this is understood of the very person of Saint John, yet even herein is understood that the Lord will have his church likewise instructed and taught by other preachers also. This pertains to the consolation of the faithful, who shall live in the days of Antichrist and the residue. Thomas of Aquin also: In Saint John himself, says he, other preachers are understood, whom the Lord in the time of Antichrist will have to preach instantly to great and small. So much says Thomas.

\index{Arethas of Caesarea}\index{Resurrection}Arethas, Bishop of Caesarea, an expositor of this book, recounts concerning this passage from Saint John that the opinion of the common people was that Saint John, along with Enoch and Elijah, would return into the world before the judgment, namely, corporally, and earnestly and constantly to preach against Antichrist. The same Arethas repeats with a more plentiful exposition, where in the eleventh chapter. Where in chapter 44 of Ecclesiastical History it is written that Enoch was translated that he might teach the Gentiles, many have interpreted it as though he should corporally return, that he might teach the Gentiles against Antichrist; whereas by the very translation made in times past he teaches rather the Gentiles that there is another life prepared for the servants of God, and that the same is also due for the bodies, since Enoch was translated both in body and soul: against the opinion of Epicurus and the mad world, supposing no other life to remain after this, and that the bodies do putrefy and never rise again. This Enoch seems to come spiritually to that last age, for that the Lord himself prophesied that a like thing should come unto it, as happened before the deluge or flood of Noah. For just as many then, being careless, scorned the judgments of God, neither feared they any peril, nor hoped for any better life, so it comes to pass also in the last age, in which Enoch constantly preaches through those who establish and maintain eternal life and the resurrection of bodies against the Epicureans.

\index{Justification}Elijah appeared in glory with our Savior Christ to three chosen Apostles on Mount Tabor; neither is it to be thought that about the end of the world he must be thrust out of the heavenly palace, and again be subject to corruption, and subjected to the cruel hands of Antichristians, which might tear him in pieces. For just as in the time of our Savior Christ Elijah in virtue and spirit, I mean Saint John the Baptist, went before Christ the Lord, so also before the judgment Elijah shall preach in those who are endowed with the spirit and virtue of Elijah, who shall draw away the minds of all men from the worship of creatures to the adoration of the eternal and only God. Elijah cried out, “How long halt you on both sides?” If the Lord be God, follow him; if Baal be God, follow him. And now the Elijahs shall cry: if Christ is the perfection of the faithful, what need is there of human inventions and constitutions to work a perfection? If Christ is our justification, satisfaction, purification, our only mediator and redeemer, why are these things attributed to human merits? Why are saints accounted intercessors in heaven? Why is salvation ascribed to many other stinking things? Elijah cried out, “How long halt you on both sides?” As though he should say, it is not lawful to part your hearts between two gods, neither is it lawful to attribute all things of life and of salvation but unto God alone. The fellowship of the kingdom is in this case envious indeed. The Elijahs shall cry: if righteousness is of the law, Christ died in vain. No man can serve two masters. Christ shall profit you nothing, which seek salvation in the traditions of men. Come to Christ: he is the perfection of the faithful, and in him we are complete. And just as Elijah grievously accused Ahab, Jezebel, and the Balamites, so shall the Elijahs most sharply inveigh against kings, bishops, idolaters, and Antichristians. Thus I say, Elijah comes again, has come, and shall come before the judgment.

\index{Apocalypse (Book of)}\index{Judgment, Last}Neither shall John the Apostle prophesy otherwise before the judgment. He shall not return to the earth in his body from heaven; but preachers endowed with John’s doctrine shall renew all his doctrine, they shall expound those things which he has left to the church, written in his Gospel, in his Epistles, and in the Apocalypse. This book has for a time lain hidden, and has been disdained by good and learned men; yet preaching what is contained and set forth in this book will bring it to light and establish it, as is plainly seen to have happened with many others in our memory. And from all these things we clearly perceive how Antichrist must be opposed and slain, not with carnal armor, but with spiritual means: namely, by sincere doctrine, framed after the example of Enoch, Helie, and John, and taken from the holy scriptures. Of this we shall speak more fully in chapter 11. Briefly, John’s doctrine about the last judgment shall be renewed and become known to the world, despite and against their will. And by John’s doctrine we understand the entire evangelical and apostolic doctrine, in the writing and setting forth of which John also exerted a singular effort among the most excellent.

\index{Papacy}And in the meantime, within the same vision, the entire manner of faithful and lawful preachers being matched with Antichrist is depicted, showing what they ought to be and with what qualities they should be equipped. First, John is called by a voice spoken to him from heaven, with a commandment to go. Therefore, God’s vocation is chiefly necessary, lest any man should take upon himself this office with an evil affection. Moses was called, the prophets and Apostles were called: some indeed immediately from God, not of men, nor by men; some of God also, but yet by men and of men. The apostles of Christ were not called of men boasting of lawful succession, from Caiaphas, Annas, and the college of priests; nevertheless, they had their vocation from Christ, and approved their vocation in deed, by preaching the truth. Therefore, although we cannot at this day refer our vocation to the Pope and bishops, boasting of lawful succession, yet forasmuch as we are able to approve it in very deed, and by the testimonies of Christ, that our doctrine is Christ’s doctrine, and therefore that our ministry is lawful, we care not a whit for their opprobrious and railing words, which cry that we are not called, that we are not ordained by the Pope.

But to those who are called, a sure commandment is given: to take the book—not every book, but the book open, and that from the hand of the angel, and again from the angel standing upon the sea and land. That angel is Christ the Lord, Lord of the whole earth, of the sea and all things contained therein. He, with his hand, offers to his ministers a book open—namely, the holy scripture, and especially his sacred and holy gospel, wrapped with no darkness, neither closed, but right manifest to those who will see. For although, because of the antiquity of the language, the propriety of speech, the figures, and the rites, places, things, and stories, some passages may appear somewhat difficult, what does this darken or obscure the mystery of faith and salvation, which is nevertheless most open and plain? Who does not understand what he should believe, what he should do, and how he should pray, even from the articles of the faith, the ten commandments, and the Lord’s Prayer? The sum of faith and of doctrine is certain and plain.

This book, therefore, opened, Christ offers to his ministers. Saint John speaks of a little book, not of a book. For if you compare the holy Bible, especially the gospel book, with other laws, books, and especially the decrees and decretals of the Pope, the little book of the holy gospel will seem very small. Primasius, expounding this passage, seems to understand the truth of the law and prophets manifested in Christ; therefore he says not now, as before, that he takes the sealed, but the open book. For Christ is the end of the law for righteousness to all who believe, and so forth. Therefore, the Lord Christ himself gives to sincere preachers no other preaching than his own, namely, the Evangelical. For he is the light and redeemer of the world, righteousness and life, and there is no other salvation. This preaching is not drawn from or taken from others than from the hands of the angel, not from the hands of the Pope or bishops. Christ says, “Go forth into the whole world and preach the gospel to every creature, teaching them to keep all things which I have commanded you.”

Now also is required obedience of the ministers, that they obey the commandment of God: and that they seek and receive that which they are commanded to ask and receive. It is in vain that some look for a drawing and working of salvation outwardly, and without effort on their part to be finished, through the only invisible operation of God. If God will have me blessed and just, they say, let him work in me what he will. Moreover, they themselves are not careful how they should apply themselves to the grace of God working by grace. It is against their ungodliness that we hear how John applies himself to the commandments of God, not without grace. For he goes to the angel and says, “Give me the book.” For the Lord must be prayed to; we must read diligently, as Paul also commands; we must learn and obey the commandments of God, and not tarry until God without us draws us.

\index{Jerome, Saint}\index{Ezekiel (Book of)}And the Lord denies nothing to those who are willing, do, and are diligent, as he says in the Gospel: “I will give you a mouth and wisdom that your adversaries will not be able to contradict.” Moreover: “My good Father will give his Holy Spirit to those who desire it from him.” Therefore, the angel now says: “Take the book.” Nevertheless, he adds another commandment: “Eat it.” He alludes to the second and third chapter of Ezekiel, where the prophet is likewise commanded by God to eat a book offered to him. For John invents nothing new here. Jerome says that to eat a book is to lay up the understanding of the scriptures in the secret bowels or entrails. He seems by a trope both to intimate an earnest desire and to signify a singular diligence. For we devour with a greedy desire such things as we have long and much coveted to eat. They are said to have devoured books and authors, which they have perfectly learned and can. We say in Dutch, “Er hat den Galen, oder Priſciane gar freſſen,” that is to say, he has learned them perfectly. It is required therefore of the preachers that they learn the holy scriptures with a desire, and that they learn and remember them whole and exactly. Without a desire and fervor of mind, you shall profit little in the study of holy scriptures; and unless you learn the Gospel exactly, you shall preach it unprofitably. The ministers therefore may be ashamed of their ignorance, which is given to idleness, taverns, hunting, dressing [?], and other worse things, than to the study of holy scriptures. They being far unlike the apostle John, shall in this warfare against Antichrist win small renown, unless they do awake out of their profane sleep, and faithfully do their duty, without doubt most holy.

\index{Repentance}Neither is the effect of the ministry and word preached concealed here. It is sweet in the mouth as honey. For David also sings: “The judgments of the Lord are to be desired above much gold and precious stones, and sweeter than honey or honeycombs.” This sweetness is always felt in the inward man, and the faithful, enlightened by the truth, have always continual joy. But we must not conceal what it seems to the flesh, and what the effect thereof is in the outward man. It truly makes one violently agitated—which is also a phrase of speech, to which others respond, signifying that what is proposed to us is both painful and grievous. The word of God therefore brings the mortification of the flesh, trials, pain, the cross, and innumerable adversities, which with a strong and constant patience we must overcome. For the Lord in the Gospel preached repentance or mortification, and among other things made much mention of persecutions with which his followers would always be exercised. Primasius: “When you have devoured the book,” he says, “you will indeed be delighted with the sweetness of the divine word, and with the hope of salvation promised, and the pleasant taste of God’s righteousness; but without doubt you will feel bitterness when you begin to preach both to the devout and the ungodly.” For the preaching of God’s judgment, once heard, doubtless through the bitterness of repentance, some being converted to better are changed; and others again, being offended, are more hardened, and bear great hatred and malice toward the preachers. The wise man says, “You shall rebuke a wise man, and he will love you; reprove a fool, and he will hate you,” and so forth.

These things are not said only, but also done and felt; for John says, “And when I had devoured it, my belly was made bitter.” And we feel at this day the most grievous hatred from powerful men, namely, of spiritual fathers and temporal princes. Many are driven into exile, innumerable are shut up in prisons, and an infinite multitude are slain with sundry kinds of deaths. All these things the prophets foretold should come to pass, our savior himself in the Gospel gave us warning thereof; the Lord here tells us again the same tale. Therefore let us be strong and steadfast in the Lord, and fight against Antichrist until the end of our life. The Lord will not forsake us, lest we should be overcome by those adversities, having told us of them diligently beforehand. And thus must they be instructed who shall war against Antichrist before the last judgment.

As I stated at the beginning of the sermon, here is a brief explanation of the vision. For the angel says to John, “You must prophesy again to the Gentiles, and so forth.” By this vision, he says, I would declare nothing else but that you must preach again to the world, first by yourself in Asia, after you shall return from exile; secondly by faithful ministers even to the world’s end, who shall spread abroad this doctrine now set forth by you, with sundry tongues through nations, and therewith shall defeat Antichrist. Those accustomed to reading the scriptures know that to prophesy is taken to mean to preach. For prophecy is preaching; they were in times past called prophets, which at this day are preachers, as we may gather from 1 Corinthians 11 and 14. And the doctrine of John is turned into the Syrian, Ethiopian, Egyptian, German, Spanish, French, English, Italian, to be short, in a manner into all languages: in all these, John preaches at this day by faithful ministers. The Gentiles, however barbarous and rude, hear John teaching, and so do the people of many nations. All these receive not a little comfort in these most dangerous days of Antichrist, and have received of them also before this time, which long sins renewed the apostolic doctrine against Antichrist. The same doctrine is brought at this day, and was brought in times past also unto kings and Popes, though they kicked and spurned against it. The thing I speak is not doubtful. For we both hear and see these things even at this day. Histories also report many things hereof. Land and glory be to God. Some copies in the Latin are corrupt, which have \textit{Igitur} for \textit{Iterum}. For John said, “You must prophesy [illegible],” which signifies \textit{Iterum} again, not \textit{Igitur}. For he signifies that he, being dead also, must preach to many nations in sundry tongues, by faithful ministers that shall fight against Antichrist. The Lord assist with his spirit all godly preachers of the Evangelical verity and apostolic doctrine. Amen.
